Automatisering gaat om data verwerking. We hebben in voorgaande lessen gezien dat Python data kan opslaan in variabelen en daar kunnen we dan bewerkingen op los laten. Ook hebben we data opgevraagd aan de gebruiker via de \texttt{input} functie. Er zijn echter nog vele andere manieren waarop we data kunnen binnenhalen binnen ons programma. We kunnen bijvoorbeeld een bestand inlezen of data aangeleverd krijgen via opties. Dit document gaat over deze manieren om data op te nemen in een programma en daarna deze data te verwerken.

